\section{Buridan’s Donkey}

Free will is not what you suppose.

\begin{quotex}
Should two courses be judged equal, then the will cannot break the deadlock, all it can do is to suspend judgement until the circumstances change, and the right course of action is clear.
\flright{\textsc{Jean Buridan}, c. 1340}
\end{quotex}


Buridan's claim led to the development of a thought experiment that goes by the name ``Buridan's Ass''\footnote{\url{https://en.wikipedia.org/wiki/Buridan\%27s\_ass}}. However, that word has an ambiguous meaning in American English so I decided to rename the experiment to ``Buridan's Donkey''. For non-native speakers, if you take your son to the petting zoo, you wouldn't say to him, ``Look at that nice ass over there!'' You should say instead, ``Look at that cute donkey over there!''

The simplest version is to suppose you place two bales of hay equidistant from a hungry donkey. Since he would have no rational basis to choose between the two bales, he would starve to death. Philosophically, this story attempts to convert the metaphysical question of free will into a science experiment. If the experiment were actually carried out, the donkey would not starve. Hence one of the premises must be false, to wit:

\begin{itemize}


\item The donkey does not have free will

\item The donkey is not rational

\item The two options are unequal
\end{itemize}


The donkey is not rational so he does not have free will. Let's look at al-Ghazali's version which I can relate to:

\begin{quotex}
Suppose two similar dates in front of a man, who has a strong desire for them but who is unable to take them both. Surely he will take one of them, through a quality in him, the nature of which is to differentiate between two similar things.
\end{quotex}


I do have a bag of Turkish dates. Whenever I desire one, I stick my hand in the bag and pull one out at random. What I should do is put two of them each equidistant from me and let you know what happens; maybe I can even record it and put it on Tik Tok. Al-Ghazali says something much more intelligent. He assumes I will take one of them through a ``quality'' that can choose the correct date. So let's investigate that quality.

\paragraph{The Intellectual Soul}


A donkey only has an animal soul, so it is incapable of making a rational judgment. It does not deliberate options and then make a choice, which is a gnomic will. Rather it acts by instinct, which is its ``quality'' that ``chooses'' the correct bale of hay.

A human being has an intellectual soul as well as an animal soul. In a well-ordered person, the animal soul is subservient to the intellectual soul. A being with intellect needs to be free because:

\begin{itemize}


\item The intellect chooses the true rather than the false

\item The intellect chooses the good rather than the bad
\end{itemize}


You can see, then, that if the intellect cannot choose between the true and the false, there is no possibility of knowledge. The most important thing to know is what is the good. It would be pointless to know the good, but to be unable to act on it.

\paragraph{Physical Causation}


The most primitive objection to free will is to assume that the brain works entirely through physical causes and executes ``algorithms''. The corollary is that the mind is merely a passive observer of those processes and can do nothing to change them. That is a horrifying thought if you understand it and try to live by it.

While it purports to be ``scientific'', there is no science behind it at the moment. There is no testable theory. More critically, there is no scientific measurement to determine if a subject is conscious or not. Nor is there an independent way to determine if an idea is true or good.

Moreover, who is executing the program?

\paragraph{Free Choice}


A free choice is one made by the conscious subject or person. The free person knows the good and acts on it. That is our real and true nature. Freedom, therefore, is to act in accordance with our true nature.

Choosing the bad, therefore, cannot be free because the choice is based on ignorance or on weakness that yields to the inclinations of the animal soul.

\emph{So you can be a person or a donkey, like Pinocchio.}

How to tame your inner donkey: \url{https://www.youtube.com/watch?v=EDK9KOfknTw}

How to make up your mind: \url{https://www.youtube.com/watch?v=C0cJ\_\_JBnWg}

\flrightit{Posted on 2021-10-10 by Cologero}

